\usepackage{mathrsfs}
\let\mathcal\mathscr

\newcommand{\Changel}{\mathcode`l="7160}
\newcommand{\Changelback}{\mathcode`l="716C}
\newcommand{\NoChangel}[1]{%
\expandafter\let\csname old\string#1\endcsname=#1
\let#1=\relax
\newcommand{#1}{\mathcode`l="716C\csname old\string#1\endcsname\mathcode`l="7160 }%
}

\NoChangel{\log}
\NoChangel{\ln}
\NoChangel{\lim}
\NoChangel{\limsup}
\NoChangel{\liminf}
\NoChangel{\varlimsup}
\NoChangel{\varliminf}
\NoChangel{\varinjlim}
\NoChangel{\varprojlim}

\newenvironment{enumeratei}
{\bgroup\def\theenumi{\roman{enumi}}\def\theenumii{\arabic{enumii}}\begin{enumerate}}
{\end{enumerate}\egroup}
\newcommand\mto{\mathchoice{\longmapsto}{\mapsto}{\mapsto}{\mapsto}}

\newcommand{\psfrac}[2]{\sfrac{(#1)}{#2}}

\DeclareMathOperator{\GL}{GL}

\usepackage{tikz}
\usetikzlibrary{calc,matrix,arrows,shapes,decorations.pathmorphing,decorations.markings,decorations.pathreplacing}

\newcommand{\RR}{\mathbf{R}}
\newcommand{\CC}{\mathbf{C}}
\newcommand{\NN}{\mathbf{N}}
\newcommand{\ZZ}{\mathbf{Z}}
\newcommand{\PP}{{\mathbb{P}}}

\DeclareMathOperator{\pgcd}{pgcd}
\DeclareMathOperator{\ord}{ord}
\DeclareMathOperator{\rot}{rot}
\DeclareMathOperator{\ind}{Ind}

\newcommand\Res{\operatorname{Res}}
\newcommand{\Resk}[1][k]{\Res^{#1}}
\newcommand{\appres}[1][g]{\mathfrak{R}_{#1}}
\newcommand{\espres}[1][g]{\mathcal{R}_{#1}}
\newcommand{\barmoduli}[1][g]{{\overline{\mathcal M}}_{\!#1}}
\newcommand{\omoduli}[1][g]{{\Omega\mathcal M}_{\!#1}}

\newcommand{\Weierstrass}{Weierstra\ss{}}

\newtheorem{thm}{Théorème}[section]
\newtheorem{cor}[thm]{Corollaire}
\newtheorem{prop}[thm]{Proposition}
\newtheorem{lem}[thm]{Lemme}

\theoremstyle{definition}
\newtheorem{defn}[thm]{Définition}
\newtheorem{rem}[thm]{Remarque}
\newtheorem{ex}[thm]{Exemple}

